\documentclass[12pt]{beamer}

\usepackage[utf8]{inputenc}
\usepackage{graphicx}
\usepackage{wrapfig}

\usetheme{Rochester}

\title{Tutorium - Einführung in die Informatik}
\author{Kilian Gärtner, Rosario Raulin}
\date{Wintersemester 2012/13}

\begin{document}

\frame{\titlepage}

\section{Wer sind wir?}

\begin{frame}

	\frametitle{Wer wir sind (1)}

	\begin{itemize}
		\item Rosario Raulin
		\item Informatik-Student im 3. Semester
		\item Studienprofil: Webgründer
		\item E-Mail: raulin@st.ovgu.de
		\item Facebook, Twitter, ...
	\end{itemize}

	\begin{flushright}
		\includegraphics[scale=0.1]{src/img/babygnu}
	\end{flushright}

\end{frame}

\begin{frame}

	\frametitle{Wer wir sind (2)}

	\begin{itemize}
		\item Kilian Gärtner
		\item Informatik-Student im 3. Semester
		\item E-Mail: kilian.gaertner@st.ovgu.de
		\item Facebook, Minestar.de, ...
	\end{itemize}

\end{frame}

\begin{frame}
	\frametitle{Agenda}
	\tableofcontents
\end{frame}


\section{Ziel und Motivation}

\begin{frame}

\frametitle{Worum geht's eigentlich?}

	\begin{quote}
		Ein Tutorium oder Tutorat ist an einer Hochschule eine Lehrveranstaltung
		im Grundstudium, in der ein fortgeschrittener Student eine
		Lehrveranstaltung unterstützt, indem er mit den Teilnehmern
		Grundkenntnisse vertieft und -fertigkeiten einübt.

		\begin{flushright}
			\scriptsize Quelle: Wikipedia
		\end{flushright}

	\end{quote}
	
\end{frame}

\begin{frame}

	\frametitle{Oder anders ausgedrückt...}

	\pause
	\begin{center}
		\includegraphics[scale=0.42]{src/img/help}
	\end{center}
	\pause

	\begin{itemize}
		\item wöchentlich, mittwochs, 11.15 - 12.45 Uhr, G29-336
		\item alle 2 Wochen "Pflicht" (wir zwingen Euch aber nicht)
	\end{itemize}

\end{frame}

\begin{frame}

	\frametitle{Unsere Ziele}

	\begin{wrapfigure}{l}{60px}
		\includegraphics[scale=0.15]{src/img/doc}
	\end{wrapfigure}

	Unsere Ziele:
	\begin{enumerate}
		\item Abbrecherquote senken
		\item sehr gute Prüfungsergebnisse
	\end{enumerate}

\end{frame}

\begin{frame}
	\frametitle{Was wir bieten}

	\begin{wrapfigure}{r}{0pt}
		\includegraphics[scale=0.55555]{src/img/tools}
	\end{wrapfigure}

	Wir bieten Euch:
	\begin{itemize}
		\item praktisch Umsetzung des Vorlesungsstoffs
		\item Hilfe in Theorie und Praxis
		\item Programmiererfahrung (Kilian, siehe minestar.de)
		\item kompetente Ansprechpartner ;-)
	\end{itemize}

\end{frame}

\begin{frame}
	\frametitle{Take-Home-Message}
	\begin{quote}
		Wenn man etwas nicht weiß, so kann man fragen; wenn man etwas nicht
		kann, so kann man es lernen.
		\newline
		\begin{flushright}
		\scriptsize Lü Buwei, chinesischer Kaufmann
		\end{flushright}
	\end{quote}

\end{frame}

\section{Vorstellungsrunde}
\section{Ausgabe der FIN-Accounts}
\section{Einführung in Eclipse}
\section{Fragen und Antworten}

\begin{frame}
	\frametitle{Vorstellungsrunde}
	\begin{enumerate}
		\item Wie heißt du?
		\item Wie alt bist du?
		\item Was studierst du?
		\item Warum studierst du das?
		\item Was sind dein Hobbys?
		\item Hast du schon Programmiererfahrung?
	\end{enumerate}
\end{frame}

\end{document}

